% Appendix B: Biographical Sketches

\chapter{Biographical Sketches}
\label{app:biographical}

This appendix provides brief biographical context for the key figures mentioned throughout the book. Entries are alphabetical by surname.

\section*{Alon, Noga (b. 1956)}
Israeli mathematician. Fundamental contributions to combinatorics, spectral graph theory, and the probabilistic method. Co-proved the existence of expanders via the Alon-Boppana bound. The Alon-Tarsi conjecture connects combinatorial nullstellensatz to graph colorings.

\section*{Arikan, Erdal (b. 1958)}
Turkish electrical engineer. Invented polar codes (2009), the first explicit construction achieving channel capacity with low-complexity decoding. The polarization phenomenon is a deep result in information theory and probability.

\section*{Babai, László (b. 1950)}
Hungarian-American computer scientist. Graph isomorphism in quasipolynomial time (2015). Co-proved MIP = NEXP, IP = PSPACE. Pioneer of interactive proofs and probabilistically checkable proofs.

\section*{Ben-Or, Michael (b. 1952)}
Israeli computer scientist. Randomized consensus breaking FLP (1983). Early work on multiparty computation and interactive proofs. Ben-Or, Goldwasser, Kilian, Wigderson (BGKW) model for multi-prover ZK.

\section*{Bennett, Charles (b. 1943)}
American physicist, IBM Research. Reversible computing, quantum teleportation, Bennett's logical depth, Maxwell's demon resolution. Bridged physics and information theory.

\section*{Berlekamp, Elwyn (1940--2019)}
American mathematician. Algebraic coding theory: Berlekamp-Massey algorithm for BCH decoding, factoring polynomials over finite fields. Co-authored \emph{Winning Ways} on combinatorial game theory.

\section*{Bonacich, Phillip (b. 1942)}
American sociologist. Eigenvector centrality (1972), Alpha centrality, power in social networks. Formalized recursive definition of influence: your importance depends on who points to you.

\section*{Bonami, Aline (b. 1947)}
French mathematician. Bonami-Beckner inequality (hypercontractivity), fundamental in analysis of Boolean functions, used in proofs of KKL theorem, noise stability, and hardness of approximation.

\section*{Brin, Sergey (b. 1973) \& Page, Larry (b. 1973)}
Stanford PhD students. Google. Revolutionized web search by treating links as votes. Built the world's most valuable company on an eigenvector.

\section*{Castro, Miguel \& Liskov, Barbara (b. 1939)}
Practical Byzantine Fault Tolerance (PBFT, 1999). First practical BFT system with $O(n^2)$ messages. Liskov: Turing Award 2008 (CLU, data abstraction, Liskov substitution principle).

\section*{Chaitin, Gregory (b. 1947)}
Argentine-American mathematician. Algorithmic information theory, Chaitin's constant $\Omega$ (halting probability), Kolmogorov complexity. ``Meta-mathematics as computation.''

\section*{Cook, Stephen (b. 1939)}
American-Canadian computer scientist. Cook-Levin theorem (1971): SAT is NP-complete. Turing Award 1982. Proved the first NP-completeness result, launching the field.

\section*{Danzig, George (1914--2005)}
American mathematician. Simplex algorithm for linear programming (1947). The simplex method is one of the top 10 algorithms of the 20th century.

\section*{Diffie, Whitfield (b. 1944) \& Hellman, Martin (b. 1945)}
Stanford. Public-key cryptography (1976): Diffie-Hellman key exchange, digital signatures, one-way functions. Turing Award 2015. ``New Directions in Cryptography.''

\section*{Dijkstra, Edsger (1930--2002)}
Dutch computer scientist. Shortest path algorithm, semaphores, structured programming, THE multiprogramming system. Turing Award 1972. ``Goto statement considered harmful.''

\section*{Dwork, Cynthia (b. 1958)}
American computer scientist. Differential privacy (2006), partial synchrony model (1988), consensus lower bounds. Turing Award 2024 (with Naor, Rothblum, Roth).

\section*{Erdős, Paul (1913--1996)}
Hungarian mathematician. Prolific collaborator (Erdős number). Probabilistic method, Ramsey theory, random graphs. ``A mathematician is a machine for turning coffee into theorems.''

\section*{Fagin, Ronald (b. 1945)}
American computer scientist. Descriptive complexity (1974): NP = $\exists$SO (existential second-order logic). Fagin's theorem connects logic and complexity. Turing Award 2018 (with S. Abiteboul, V. Vianu).

\section*{Flajolet, Philippe (1948--2011)}
French computer scientist, INRIA. Analytic combinatorics (generating functions, complex analysis). Probabilistic counting (PCSA, LogLog, HyperLogLog). Applied deep mathematics to practical algorithms. Died months after HyperLogLog publication.

\section*{Gallager, Robert (b. 1931)}
American information theorist. LDPC codes (1963), information theory textbook, queueing theory. Forgotten for 30 years; LDPC rediscovered in 1990s, now in 5G/Wi-Fi.

\section*{Garey, Michael \& Johnson, David}
\emph{Computers and Intractability} (1979) — the bible of NP-completeness. 300+ problems. Johnson: approximation algorithms, NP-completeness column in J. Algorithms.

\section*{Gentry, Craig (b. 1973)}
Stanford PhD under Boneh. First fully homomorphic encryption (2009) from ideal lattices. Bootstrapping technique. MacArthur Fellow 2010.

\section*{Goldreich, Oded (b. 1957)}
Israeli computer scientist. Foundations of cryptography (two volumes). GMW: NP $\subseteq$ CZK, parallel repetition, zero-knowledge, pseudorandomness. Israel Prize 2004.

\section*{Goldwasser, Shafi (b. 1958) \& Micali, Silvio (b. 1954) \& Rackoff, Charles}
Goldwasser-Micali-Rackoff (1985): Interactive proofs, zero-knowledge, knowledge complexity. Probabilistic encryption (GM 1982). Turing Award 2012 (Goldwasser, Micali).

\section*{Gödel, Kurt (1906--1978)}
Austrian logician. Incompleteness theorems (1931): any consistent formal system containing arithmetic is incomplete. Completeness theorem (1929). Relative consistency of AC and CH.

\section*{Groth, Jens (b. 1975)}
Danish cryptographer. Short pairing-based ZK (Groth10), zk-SNARKs (Groth16). Non-interactive ZK with constant-size proofs. Used in Zcash, Ethereum.

\section*{Hamming, Richard (1915--1998)}
American mathematician. Hamming codes, Hamming distance, Hamming bound, error-correcting codes. ``You and your research'' lecture. ACM A.M. Turing Award 1968.

\section*{Hartmanis, Juris (1928--2022) \& Stearns, Richard}
Defined time/space complexity classes (1965). P, NP, EXP, PSPACE hierarchy. Turing Award 1993. ``On the computational complexity of algorithms.''

\section*{Håstad, Johan (b. 1960)}
Swedish computer scientist. Optimal inapproximability (1998): Max-3SAT 7/8, Clique $n^{1-\epsilon}$, Set Cover $\ln n$. Switching lemma for AC$^0$ circuits. G\"odel Prize 2011.

\section*{Herlihy, Maurice (b. 1954)}
American computer scientist. Wait-free hierarchy (1991): consensus numbers of objects (registers=1, CAS=$\infty$). Herlihy-Shavit topological characterization of wait-free computability. Dijkstra Prize 2003.

\section*{Hilbert, David (1862--1943)}
German mathematician. 23 problems (1900): consistency, Entscheidungsproblem, continuum hypothesis. Formalist program. Gödel's theorems ended the program but created modern logic.

\section*{Karp, Richard (b. 1935)}
American computer scientist. 21 NP-complete problems (1972). Karp-Miller algorithm, Karp-Rabin string matching. Turing Award 1985. ``Reducibility among combinatorial problems.''

\section*{Katz, Leo (1915--1982)}
American sociologist. Katz centrality (1953): attenuated walks on graphs. Status index. Precursor to eigenvector centrality.

\section*{Kleinberg, Jon (b. 1971)}
Cornell. HITS (Hubs and Authorities, 1998), simultaneous with web search ranking algorithms. Small-world navigation, hiring algorithms, network science. Nevanlinna Prize 2006.

\section*{Knuth, Donald (b. 1938)}
American computer scientist. \emph{The Art of Computer Programming} (TAOCP). TeX, Metafont. Analysis of algorithms, literate programming. Turing Award 1974.

\section*{Kolmogorov, Andrey (1903--1987)}
Soviet mathematician. Kolmogorov complexity (1965), probability axioms, turbulence, ergodic theory. Student of Khinchin. Teacher of Levin, Arnold, Sinai.

\section*{Kuhn, Thomas (1922--1996)}
American philosopher. \emph{The Structure of Scientific Revolutions} (1962). Paradigm shifts, normal science, incommensurability. Changed how we understand scientific progress.

\section*{Lamport, Leslie (b. 1941)}
American computer scientist. Logical clocks (1978), Byzantine Generals (1982), Paxos (1998), TLA+, LaTeX. Turing Award 2013. Distributed systems pioneer.

\section*{Ladner, Richard (b. 1943)}
American computer scientist. NP-intermediate problems (1975): if P $\neq$ NP, there exist problems neither in P nor NP-complete. Gap theorem, diagonalization against reductions.

\section*{Levin, Leonid (b. 1948)}
Soviet mathematician, Kolmogorov's student. Universal search problems (1973), NP-completeness independently of Cook. Average-case complexity, one-way functions, Knuth Prize 2012.

\section*{Morris, Robert (1932--2011)}
American cryptographer. Approximate counting (1978): probabilistic counter using $O(\log \log n)$ bits. NSA cryptanalysis, Unix crypt. Father of Robert Tappan Morris (Internet Worm).

\section*{Nakamoto, Satoshi (pseudonym)}
Bitcoin (2008): Nakamoto consensus, proof-of-work, blockchain. True identity unknown. Created the first decentralized digital currency.

\section*{Ongaro, Diego \& Ousterhout, John (b. 1954)}
Raft consensus (2014): understandable Paxos. Ousterhout: Tcl/Tk, Sprite OS, log-structured file system. Raft became the standard for teaching and implementation.

\section*{Page, Larry \& Brin, Sergey}
See Brin \& Page.

\section*{Rabin, Michael (b. 1931)}
Israeli-American computer scientist. Randomized algorithms (primality testing, Byzantine agreement), non-deterministic automata, Rabin-Karp string matching. Turing Award 1976.

\section*{Rice, Henry Gordon (1920--2003)}
American mathematician. Rice's Theorem (1953): all non-trivial semantic properties of programs are undecidable. PhD under Church at Princeton.

\section*{Rivest, Ronald (b. 1947) \& Shamir, Adi (b. 1952) \& Adleman, Leonard (b. 1945)}
MIT. RSA (1977): first practical public-key encryption and signatures based on factoring. Rivest: MD5, RC4, RC5, RSA key size recommendations. Shamir: secret sharing, differential cryptanalysis, ring signatures. Adleman: DNA computing, number theory. Turing Award 2002.

\section*{Shannon, Claude (1916--2001)}
American mathematician, Bell Labs. Information theory (1948): entropy, channel capacity, source coding, noisy channel coding theorem. Cryptography (1949): perfect secrecy. Switching circuits (1937): Boolean algebra for digital circuits. ``The father of the Information Age.''

\section*{Shamir, Adi}
See Rivest-Shamir-Adleman.

\section*{Simon, Herbert (1916--2001)}
American polymath. Bounded rationality, satisficing, BACON/DENDRAL scientific discovery systems. Turing Award 1975 (with Newell). Nobel Economics 1978. ``The sciences of the artificial.''

\section*{Solomonoff, Ray (1926--2009)}
American mathematician. Algorithmic probability (1960), universal prior, inductive inference. Founder of algorithmic information theory (with Kolmogorov, Chaitin).

\section*{Stearns, Richard}
See Hartmanis-Stearns.

\section*{Turing, Alan (1912--1954)}
British mathematician. Turing machines (1936), halting problem, universal machine, Church-Turing thesis. Bletchley Park (Enigma). Morphogenesis (1952). ``Computing Machinery and Intelligence'' (1950). Turing Award named for him.

\section*{Viterbi, Andrew (b. 1935)}
American electrical engineer. Viterbi algorithm (1967): dynamic programming for convolutional codes. Hidden Markov models, speech recognition. Co-founded Qualcomm. National Medal of Science 2007.

\section*{Von Neumann, John (1903--1957)}
Hungarian-American polymath. EDVAC report (1945): stored-program architecture. Game theory (minimax), cellular automata, self-reproducing automata, quantum mechanics foundations. ``The last of the great mathematicians.''

\section*{Wallace, Christopher (1933--2004)}
Australian physicist. Minimum Message Length (MML) principle, Bayesian inference, inductive inference. MML $\approx$ Kolmogorov complexity for statistical models.

\section*{Wallas, Graham (1858--1932)}
English social psychologist. \emph{The Art of Thought} (1926): four stages of creativity — Preparation, Incubation, Illumination, Verification.

\section*{Wigderson, Avi (b. 1956)}
Israeli-American computer scientist. Hardness vs. randomness, IP = PSPACE, zero-knowledge, algebraic complexity, expander graphs. Turing Award 2023. ``Mathematics and Computation'' (2019).

\section*{Ziv, Jacob (1931--2023) \& Lempel, Abraham (1936--2007)}
Israeli computer scientists. LZ77, LZ78 compression algorithms. Dictionary-based compression. LZW (Welch) in GIF, PNG, PDF. IEEE Milestone 2004.

%======================================================================
% INDEX ENTRIES
%======================================================================
\index{biographical sketches}
\index{key figures}
\index{computer scientists}
\index{mathematicians}